APR 05 - FORMULAS/EXERCICIOS
CONVOLUCAO 2D:
$O=\lfloor\frac{I+2P-K}{S}\rfloor+1$.
Por dimensao: altura e largura separadamente.
I=tamanho input, K=kernel, P=padding, S=stride.
EX: input 6x6, kernel 3x3, P=0, S=1
O=(6-3)/1+1=4 => mapa 4x4.
EX: 32x32, K=5, P=2, S=1 => 32x32.
EX: 28x28, K=3, P=0, S=2
floor((28-3)/2)+1=13.
PARAMETROS CONV:
$(K_hK_wC_{in}+1)C_{out}$ (com bias).
EX: 3x3, 1 canal, 8 filtros => (9+1)*8=80.
EX: 3x3, 16 canais, 32 filtros:
(144+1)*32=4640.
MAXPOOL 2x2 stride 2: 8x8 -> 4x4;
recebe [1,4;2,3] => max=4.
EX CONVOLUCAO SIMPLES:
janela [1,2;3,4], kernel [1,0;0,-1]:
soma =1-4=-3 (cross-correlation usual).
DIAGRAMA
Imagem -> Conv -> ReLU -> Pool
        -> Conv -> ReLU -> Pool
        -> Flatten/GlobalAvgPool -> Dense -> classe
CICLO TREINO
Batch imagens -> forward conv/pool/dense
 -> loss -> backprop filtros e dense
 -> update -> repetir.